\chapter{Running the examples on a computer}\label{app:practical}
None of the proofs in this book needs a computer, but a computer can help in two ways: it can carry out a finite experiment exactly, and it can check a formal proof. This appendix explains the small part of the programming language Python that the book's experiments use, including exact arithmetic with fractions (Sections~\ref{appD:sec:python} and~\ref{appD:sec:fractions}); gives three complete programs, for Chapters~\ref{ch:signs}, \ref{ch:matching} and~\ref{ch:capstone} (Section~\ref{appD:sec:experiments}); and shows how to check the Lean examples of Chapters~\ref{ch:lean1} and~\ref{ch:lean2} (Section~\ref{appD:sec:lean}).
\section{The smallest amount of Python needed}\label{appD:sec:python}
The programs in this appendix are written in Python, a free programming language that runs on all common operating systems; any version from Python~3.10 on will do. The programs use only Python's \emph{standard library}, the collection of tools that comes with every Python installation, so nothing else has to be installed. This section explains just enough of the language to read the programs and run them.

\subsection*{Running a program}
A Python program is a plain text file whose name ends in \code{.py}, for example \code{signs.py}. Write it in a plain text editor; a word processor adds formatting that Python cannot read. To run the program, open a \emph{terminal}, a window for typing commands to the computer. Move to the folder that contains the file by typing \code{cd} followed by the location of that folder; for example, \code{cd Documents/math} moves into the folder \code{math} inside the folder \code{Documents}. Then type
\par\noindent\begin{minipage}{\linewidth}
\begin{lstlisting}
python3 signs.py
\end{lstlisting}
\end{minipage}\par
On some systems the command is \code{python} instead of \code{python3}. The command \code{python3 -{}-version} shows which version of Python is installed.

Short calculations can also be typed directly into Python. Type \code{python3} in the terminal with no file name. Python then displays its \emph{prompt}, usually \code{>{}>{}>}, the short sign that a program shows when it is waiting for you to type. Python carries out each line as soon as you press Enter; typing \code{2 + 2} there displays \code{4}. Type \code{exit()} to return to the terminal. The terminal has a prompt of its own, and the two must be kept apart. The terminal's prompt expects commands that start programs, such as \code{python3 signs.py}, while Python's prompt expects lines of Python. A line typed at the wrong prompt produces an error message that has nothing to do with the mathematics you wanted to test. The characters \code{>{}>{}>} belong to the prompt, so never type them yourself or copy them into a program file.

\subsection*{Names, numbers and lists}
The line \code{x = 5} stores the value $5$ under the name \code{x}. Later lines can use \code{x}, and a later line \code{x = x + 1} replaces the stored value by $6$. So the sign \code{=} in Python is an instruction to store a value. It is not a claim that two things are equal; read as an equation, $x=x+1$ would be false for every number. To \emph{test} whether two values are equal, Python uses a double sign \code{==}. The test \code{2 + 2 == 4} has the value \code{True}, and \code{2 + 2 == 5} has the value \code{False}. The other comparisons are written \code{<}, \code{<=}, \code{>}, \code{>=}, and \code{!=} for ``is not equal to.''

Arithmetic uses \code{+}, \code{-}, \code{*} for multiplication, and \code{/} for division. Two stars denote a power, so \code{2 ** 10} is $2^{10}=1024$. The function \code{abs} gives the absolute value: \code{abs(-3)} is $3$. The sign \code{\%} gives the remainder of Section~\ref{ch05:sec:euclid}: \code{17 \% 5} is $17\bmod 5=2$, and \code{-7 \% 5} is $3$, because $-7=(-2)\cdot5+3$. As in that section, the remainder on division by a positive number is never negative. The command \code{print} displays values on the screen, separated by spaces: \code{print(x, x + 1)} displays \code{6 7} when \code{x} holds $6$.

A \emph{list} is an ordered collection of values, written in square brackets. In \code{a = [1, 1, 1, -1]}, the first entry is \code{a[0]}, not \code{a[1]}: Python numbers the positions of a list from zero, just as Chapter~\ref{ch:signs} numbers the entries of a sign sequence. So \code{a[3]} is the last entry, $-1$. The \emph{length} \code{len(a)}, the number of entries, is $4$. The programs below use two operations that change a list. The line \code{a.append(9)} adds the entry $9$ at the end of \code{a}, and \code{a.pop()} removes the last entry of \code{a} and hands it back as its value.

\subsection*{Loops}
A \code{for} loop repeats a block of lines once for each value in a range. The block consists of the indented lines below the line that begins with \code{for}; Python uses this indentation, not brackets, to tell where the block ends. For example:
\par\noindent\begin{minipage}{\linewidth}
\begin{lstlisting}
a = [1, 1, 1, -1]
value = 0
for j in range(3):
    value = value + a[j] * a[j + 1]
print(value)
\end{lstlisting}
\end{minipage}\par
Here \code{range(3)} produces the values $0,1,2$: it starts at zero and stops just before $3$. So the indented line runs three times, with $j=0$, $j=1$ and $j=2$, and each time it adds one product to the running total stored in \code{value}. The products are $a_0a_1=1$, $a_1a_2=1$ and $a_2a_3=-1$, so \code{value} goes from $0$ to $1$, then $2$, then $1$, and the program prints $1$. This is the aperiodic correlation $A_1$ of $(1,1,1,-1)$ computed in Section~\ref{ch14:sec:aperiodic}.

With two numbers, \code{range(1, 14)} produces $1,2,\ldots,13$: it starts at the first number and again stops just before the second. The values of a loop can also come from a list: a block that begins with \code{for x in [2, 3, 5]:} is carried out three times, with \code{x} equal to $2$, then $3$, then $5$.

\subsection*{Functions}
A \emph{function} packages a calculation so that it can be used again with different inputs.
\par\noindent\begin{minipage}{\linewidth}
\begin{lstlisting}
def correlation(a, shift):
    value = 0
    for j in range(len(a) - shift):
        value = value + a[j] * a[j + shift]
    return value
\end{lstlisting}
\end{minipage}\par
The line beginning with \code{def} gives the function a name, \code{correlation}, and names its two inputs, \code{a} and \code{shift}. The indented lines below it carry out the calculation, and the line beginning with \code{return} gives the result. For a list \code{a} of length $n$ and a shift $s$, the loop runs over \code{range(len(a) - shift)}, which here is \code{range(n - s)} and stops just before $n-s$. So $j$ takes the values $0,1,\ldots,n-1-s$, exactly the summation range in the definition of $A_s$ (Definition~\ref{ch14:def:aperiodic}). For example, \code{correlation([1, 1, 1, -1], 1)} is $1$ and \code{correlation([1, 1, 1, -1], 3)} is $-1$, the values of $A_1$ and $A_3$ found in Section~\ref{ch14:sec:aperiodic}. The result is the correlation $A_s$ only when \code{a} is a sign sequence and $1\le s\le n-1$, the shifts allowed in the definition. The function does not check this. Given a list of $n$ signs and the shift $0$, for instance, it simply returns $n$, although the definition leaves shift zero out. Section~\ref{appD:sec:experiments} uses the same calculation, written more compactly, under the name \code{aperiodic}, next to a separate function \code{periodic} for the periodic correlation of Chapter~\ref{ch:signs}.

\subsection*{Decisions and other loops}
A line \code{if} \emph{condition}\code{:} is followed by an indented block, which is carried out only when the condition is true. Conditions can be combined with the words \code{and}, \code{or} and \code{not}, which have the meanings of Section~\ref{ch02:sec:connectives}. A \code{while} loop repeats its block as long as its condition is true. After the following lines, for example, \code{n} holds $8$, the smallest natural number whose square is at least $50$:
\par\noindent\begin{minipage}{\linewidth}
\begin{lstlisting}
n = 0
while n * n < 50:
    n = n + 1
\end{lstlisting}
\end{minipage}\par
Inside a loop, \code{break} leaves the innermost loop containing it at once, and \code{continue} skips the rest of the block and goes on with the next pass. The special value \code{None} means ``no value here.'' The test \code{x is None} asks whether \code{x} holds this special value, and \code{x is not None} asks the opposite.

\subsection*{Tuples, sets and dictionaries}
A \emph{tuple} is like a list but is written with round brackets, as in \code{(1, -1, 1)}, and cannot be changed after it has been made. Its entries are numbered from zero, as in a list, so a tuple is a natural way to store a sign sequence $(a_0,a_1,\ldots,a_{n-1})$. A function that has to give back several values at once can return them as a tuple. A piece of text is written between quotation marks, as in \code{"shortage"}. When such text is displayed as an entry of a tuple, Python puts it between single quotes, as in \code{('shortage', 1)}.

A \emph{set} records membership only, as in Section~\ref{ch03:sec:sets}: order and repetition do not matter. The set containing $1$ and $2$ is written \code{\{1, 2\}}, and \code{\{1, 2, 2\}} is the same set. The test \code{2 in s} is \code{True} when $2$ is a member of the set \code{s}, and \code{not in} tests the opposite. The union $s\cup t$ is written \code{s | t}, so \code{\{1, 2\} | \{2, 3\}} is \code{\{1, 2, 3\}}. The empty set is written \code{set()}, because empty braces \code{\{\}} mean something else in Python: an empty dictionary, described below. A loop \code{for x in s:} runs through the members of a set \code{s}, but since a set has no order, Python does not promise any particular order of visits. When the order matters, loop over \code{sorted(s)} instead. This is a list of the members of \code{s} in increasing order; for example, \code{sorted(\{3, 1, 2\})} is \code{[1, 2, 3]}.

A \emph{dictionary} is a table that pairs each \emph{key} with a value. After \code{owner = \{\}} creates an empty table, the line \code{owner[3] = 0} records the pair ``key $3$, value $0$,'' and afterwards \code{owner[3]} looks the value up. The test \code{3 in owner} asks whether the key $3$ has an entry. Asking for \code{owner[5]} when the key $5$ has no entry stops the program with an error, whereas \code{owner.get(5)} simply gives \code{None}. Python prints the table as \code{\{3: 0\}}, listing each key, a colon, and its value. The official Python tutorial explains lists, tuples, sets and dictionaries in more detail~\cite{pystructures}.

\subsection*{Adding up and checking many cases at once}
Python can write a sum much as mathematics does. The \code{for} loop that computed $A_1$ for \code{a = [1, 1, 1, -1]} can be replaced by the single expression
\par\noindent\begin{minipage}{\linewidth}
\begin{lstlisting}
sum(a[j] * a[j + 1] for j in range(3))
\end{lstlisting}
\end{minipage}\par
which reads like $\sum_{j=0}^{2}a_ja_{j+1}$ and has the same value, $1$. In the same way, \code{all(}\emph{test} \code{for s in range(1, n))} is \code{True} when the test holds for every $s$ with $1\le s\le n-1$, and \code{False} as soon as it fails for one of them. For example, with the function \code{correlation} above, \code{all(abs(correlation(a, s)) <= 1 for s in range(1, 4))} is \code{True} for \code{a = [1, 1, 1, -1]}, and \code{False} for \code{a = [1, 1, -1, -1]}, whose correlation at shift two is $A_2=-2$ (Section~\ref{ch14:sec:aperiodic}). When the range contains no values at all, as \code{range(1, 1)} does, the result is \code{True}: a statement about every member of an empty collection is vacuously true (Section~\ref{ch02:sec:contrapositive}). Finally, \code{sum(1 for x in xs if }\emph{test}\code{)} adds $1$ for each member \code{x} of the list \code{xs} that passes the test, so it counts those members. For instance, \code{sum(1 for x in [3, 8, 5, 10] if x > 4)} is $3$, because exactly three of the four numbers exceed $4$.

Two more conventions appear in the programs below. Python ignores everything on a line after the sign \code{\#}, so this sign introduces a comment for human readers. And a line \code{from fractions import Fraction} makes one tool from the standard library available to the program, here the tool \code{Fraction} for exact rational arithmetic, which Section~\ref{appD:sec:fractions} describes.

\section{An exact fraction is not a rounded decimal}\label{appD:sec:fractions}
When a program contains a number such as \code{0.1}, Python does not store one tenth exactly. It stores a \emph{floating-point number}: a binary approximation that is correct to about sixteen significant digits. The two lines
\par\noindent\begin{minipage}{\linewidth}
\begin{lstlisting}
print(0.1 + 0.2)
print(0.1 + 0.2 == 0.3)
\end{lstlisting}
\end{minipage}\par
display \code{0.30000000000000004} and \code{False}. So in an experiment that tests an equation or an inequality, a tiny rounding error can make a true statement look false, or a false one look true.

When the numbers involved are rational, the standard library offers exact arithmetic instead~\cite{pyfractions}. A \code{Fraction} stores an integer numerator and a positive integer denominator, reduced to lowest terms, and it computes by the usual rules for fractions, with no rounding at any stage. For example, \code{Fraction(1, 3)} is the fraction $1/3$, and \code{Fraction(0)}, with the denominator $1$ left out, is $0/1$, that is, zero. Python's integers are exact as well, however large they become.
\par\noindent\begin{minipage}{\linewidth}
\begin{lstlisting}
from fractions import Fraction
print(Fraction(1, 10) + Fraction(2, 10) == Fraction(3, 10))
x = Fraction(1, 3)
y = (x * x + 1) / 3
print(y)
\end{lstlisting}
\end{minipage}\par
This program displays \code{True} and then \code{10/27}. The last three lines apply the map $T(x)=(x^2+1)/3$ of Section~\ref{ch12:sec:example} to the iterate $x_1=1/3$, so \code{y} is $x_2=T(x_1)$, and Python finds exactly the value $10/27$ computed by hand in that section. (Typed at Python's prompt without \code{print}, the same value is shown as \code{Fraction(10, 27)}.)

Two details matter. First, dividing one ordinary integer by another with \code{/} gives a floating-point number, so \code{1 / 3} is \code{0.3333333333333333}. But a \code{Fraction} combined with an integer or with another \code{Fraction} gives again a \code{Fraction}, which is why the line \code{y = (x * x + 1) / 3} can use the plain integers $1$ and $3$ and still be exact. A \code{Fraction} combined with a floating-point number such as \code{0.5}, on the other hand, gives a floating-point number, so do not mix the two kinds. Second, one tenth must be written \code{Fraction(1, 10)}. Writing \code{Fraction(0.1)} instead converts the floating-point number \code{0.1} into a fraction exactly as it is stored, and \code{print(Fraction(0.1))} displays
\par\noindent\begin{minipage}{\linewidth}
\begin{lstlisting}
3602879701896397/36028797018963968
\end{lstlisting}
\end{minipage}\par
The denominator is $2^{55}$, and this fraction is close to one tenth but not equal to it.

The experiments of Section~\ref{appD:sec:experiments} keep every comparison exact: the sign-sequence counts use only integers, the iteration uses fractions, and decimals appear only when a result is displayed.

\section{Three experiments}\label{appD:sec:experiments}
Each of the following programs is complete: save it in a file of its own and run it as described in Section~\ref{appD:sec:python}.

\subsection*{Counting sign sequences}
The following program produces the table of Section~\ref{ch14:sec:experiment}.
\begin{lstlisting}
from itertools import product

def aperiodic(a, s):
    n = len(a)
    return sum(a[j] * a[j + s] for j in range(n - s))

def periodic(a, s):
    n = len(a)
    return sum(a[j] * a[(j + s) % n] for j in range(n))

def is_barker(a):
    return all(abs(aperiodic(a, s)) <= 1 for s in range(1, len(a)))

def is_circulant_row(a):
    return all(periodic(a, s) == 0 for s in range(1, len(a)))

for n in range(1, 14):
    seqs = list(product([1, -1], repeat=n))
    barker = sum(1 for a in seqs if is_barker(a))
    circ = sum(1 for a in seqs if is_circulant_row(a))
    print(n, barker, circ)
\end{lstlisting}
The first line makes the tool \code{product}, from the part of the standard library called \code{itertools}, available to the program. Called as \code{product([1, -1], repeat=n)}, it fills $n$ positions with entries of the list \code{[1, -1]} in every possible way, so it produces all $2^n$ sign sequences of length $n$, one at a time, each as a tuple; for $n=2$ these are \code{(1, 1)}, \code{(1, -1)}, \code{(-1, 1)} and \code{(-1, -1)}. The line \code{seqs = list(...)} collects them into one list.

Exact arithmetic cannot show that a program computes the quantity the text defines, or that its loops run over exactly the intended objects. So before trusting the output, read each function against the definition it is meant to carry out.
\begin{itemize}
\item The function \code{aperiodic} is the function \code{correlation} of Section~\ref{appD:sec:python}, written more compactly with \code{sum}; as checked there, $j$ runs through the summation range of $A_s$ in Definition~\ref{ch14:def:aperiodic}.
\item In \code{periodic}, the index $j$ runs through all of $0,1,\ldots,n-1$, and \code{(j + s) \% n} is the position $(j+s)\bmod n$ of Definition~\ref{ch14:def:periodic}. This is where the wrap-around happens.
\item The test \code{is\_barker(a)} is \code{True} exactly when $\abs{A_s}\le1$ for every shift $s$ with $1\le s\le n-1$, the shifts that \code{range(1, len(a))} produces; this is Definition~\ref{ch14:def:barker}.
\item Likewise, \code{is\_circulant\_row(a)} is \code{True} exactly when $C_s=0$ for every shift $s$ with $1\le s\le n-1$. By Section~\ref{ch14:sec:circulant}, these are exactly the sequences that form the first row of a circulant Hadamard matrix, hence the name.
\end{itemize}
For each length $n$ from $1$ to $13$, the main loop counts in \code{barker} and \code{circ} the sequences that pass each test and prints one line with $n$ and the two counts; the thirteen lines agree with the rows of the table. The first line is \code{1 2 2}: for $n=1$ there are no shifts to test, so \code{all(...)} is vacuously \code{True}, and both sequences of length one pass both tests, just as in Chapter~\ref{ch:signs}. The fourth line is \code{4 8 8}: of the $16$ sign sequences of length four, eight are Barker sequences, and eight have $C_1=C_2=C_3=0$.

To go further, raise the limit \code{14} in the line \code{for n in range(1, 14):}. Each additional length doubles the number of sequences, each a little longer to check, so the running time slightly more than doubles and the memory doubles, since \code{seqs} holds all $2^n$ sequences of one length at once. With the limit \code{21} (lengths up to $20$, about a million sequences) the program finishes well within a minute on an ordinary computer, but at length $25$ the list alone has more than $33$ million entries and needs several gigabytes.

\subsection*{An assignment search that returns a certificate}
The next function carries out the procedure of Section~\ref{ch08:sec:search}. Students are numbered $0,1,2,\ldots$, and \code{options[s]} is the set $N(s)$ of tasks acceptable to student \code{s}.
\begin{lstlisting}
def assign(options):
    # options[s] is the set of tasks acceptable to student s.
    # Returns ("assignment", task_of) or ("shortage", A, N(A)).
    owner = {}     # task -> student who currently holds it
    task_of = {}   # student -> task that student currently holds
    for root in range(len(options)):
        came_from = {}        # task -> student from which it was reached
        reached = [root]      # students reached by the search
        to_explore = [root]   # reached students not yet inspected
        free_task = None
        while to_explore and free_task is None:
            student = to_explore.pop()
            for task in sorted(options[student]):
                if task in came_from or task_of.get(student) == task:
                    continue
                came_from[task] = student
                if task not in owner:
                    free_task = task
                    break
                reached.append(owner[task])
                to_explore.append(owner[task])
        if free_task is None:     # the search failed
            needed = set()        # will become N(A)
            for s in reached:
                needed = needed | options[s]
            return ("shortage", reached, needed)
        task = free_task          # flip the augmenting path back to the root
        while task is not None:
            student = came_from[task]
            old_task = task_of.get(student)
            owner[task] = student
            task_of[student] = task
            task = old_task
    return ("assignment", task_of)

print(assign([{1, 2}, {1}, {2, 3}]))
print(assign([{1, 2}, {1}, {2}]))
\end{lstlisting}
\emph{The matching.} Two dictionaries describe the current matching $M$: \code{owner[t]} is the student who holds task \code{t}, and \code{task\_of[s]} is the task that student \code{s} holds. A task without an entry in \code{owner} is unused, and a student without an entry in \code{task\_of} is unmatched. Both start empty, so $M$ starts as the empty matching. The students then take turns as the root of a search, in the order $0,1,2,\ldots$, and each is still unmatched when its turn comes, as the procedure requires: a flip changes the task of students who already hold one, and gives a first task only to the root.

\emph{The search.} The dictionary \code{came\_from} records each reached task with its predecessor, the student from which it was reached; the list \code{reached} holds the reached students, and \code{to\_explore} those not yet inspected. A list used as a condition counts as true when it is not empty, so the \code{while} loop runs as long as some reached student awaits inspection and no unused task has been found. Each pass takes a student from \code{to\_explore} and inspects it, taking its acceptable tasks in increasing order thanks to \code{sorted}, as in the trace of Section~\ref{sec:matching-trace}. The \code{continue} line skips a task that has already been reached, and also the task the student holds, since the edge to that task is matched.

Every other task is marked as reached, with the current student as its predecessor. If the task is unused, the search has succeeded: \code{break} ends the inspection, and the \code{while} loop stops too, because \code{free\_task} is no longer \code{None}. If the task is used, its owner is marked as reached and waits to be inspected; by the argument of Section~\ref{ch08:sec:search}, this owner has not been reached before, so \code{reached} never lists a student twice. Its predecessor need not be recorded separately: it is the task the owner holds, which \code{task\_of} already records. Because \code{pop()} takes the most recently reached of the waiting students, this is not breadth-first search, but as Section~\ref{ch08:sec:search} points out, the order of inspection does not matter.

\emph{The flip.} If an unused task was found, the last loop flips $M$ along the augmenting path, walking backward from the unused task to the root. At each step, \code{student} is the predecessor of the current task; this student gives up its old task, saved as \code{old\_task}, and takes the current one. The walk continues from \code{old\_task}, the task from which this student was reached, and in the next pass \code{owner} records the new holder of that task, its predecessor, so no out-of-date entry is left behind. At the root, \code{old\_task} is \code{None}, because the root held no task, and the loop stops. The two dictionaries now describe the flipped matching, which by Lemma~\ref{thm:augment} has one more edge than $M$.

\emph{A failed search.} The \code{while} loop stops only when an unused task has been found or \code{to\_explore} is empty, that is, when every reached student has been inspected; so the search is complete in the sense of Section~\ref{ch08:sec:search}. Suppose that it ends without an unused task, and let $A$ be the set of reached students. The three lines that build \code{needed} collect every task accepted by some student of $A$, so \code{needed} is $N(A)$. By Proposition~\ref{ch08:prop:shortage}, $N(A)$ has exactly one member fewer than $A$, so $A$ is a shortage set. The function returns $A$ and $N(A)$ at once, without trying the remaining students, because a shortage set already proves that no matching gives every student a task (Section~\ref{ch08:sec:shortage}).

The two test lines use the examples of Section~\ref{sec:matching-trace}, with the students $a,b,c$ numbered $0,1,2$ and the tasks keeping their names $1,2,3$. The word in quotation marks at the start of each answer names its kind. The program displays
\par\noindent\begin{minipage}{\linewidth}
\begin{lstlisting}
('assignment', {0: 2, 1: 1, 2: 3})
('shortage', [2, 0, 1], {1, 2})
\end{lstlisting}
\end{minipage}\par
In the first line, \code{\{0: 2, 1: 1, 2: 3\}} pairs student~$0$, that is $a$, with task~$2$, then $b$ with task~$1$, and $c$ with task~$3$: the assignment of Section~\ref{sec:matching-trace}, reached by a different route, because the program takes the roots in the order $a,b,c$ rather than $a,c,b$. The search from $a$ finds the unused task~$1$ at once, and $a$ takes it. The search from $b$ reaches task~$1$ from $b$, then its owner $a$, and then the unused task~$2$ from $a$, so it ends with \code{came\_from} equal to \code{\{1: 1, 2: 0\}}. The flip gives task~$2$ to its predecessor $a$, whose old task is~$1$, and then task~$1$ to its predecessor $b$, whose old task is \code{None}, so the walk stops. Finally, the search from $c$ reaches task~$2$, now held by $a$, and then the unused task~$3$, which $c$ takes.

In the second test, task~$3$ has been removed from $c$'s options. The searches from $a$ and $b$ run as before and give the matching $\{a\text{--}2,\ b\text{--}1\}$. The search from $c$ reaches task~$2$ and its owner $a$, then $a$'s other task~$1$ and its owner $b$, who accepts no further task, so the search fails. The list \code{[2, 0, 1]} contains the reached students $c$, $a$, $b$ in the order of reaching, and \code{\{1, 2\}} is the set of tasks they accept: three students, but only two tasks. This is the shortage set of Section~\ref{sec:matching-trace}, although the program reached it from a different matching.

Neither kind of answer has to be taken on trust, because each is a certificate in the sense of Section~\ref{ch08:sec:shortage}. To check an assignment, make sure that every student has a task that the student accepts and that no task appears twice. To check a shortage, count the tasks that the reported students accept: there must be fewer tasks than students.

\subsection*{The biased iteration with exact fractions}
The last program checks the numerical experiment of Section~\ref{ch18:sec:experiment}. It computes the iterates of the biased update
\[
y_{n+1}=\frac{3y_n+1}{4}+\frac1{100},\qquad y_0=0,
\]
with exact fractions. For each $n$ from $0$ to $20$, it checks the closed form $y_n=\frac{26}{25}\bigl(1-(\frac34)^n\bigr)$ and the error certificate, the bound of part~2 of Theorem~\ref{thm:inexact}, which here reads
\[
\abs{y_n-1}\le\Bigl(\frac34\Bigr)^n e_0+\eta\,\frac{1-(3/4)^n}{1-3/4},
\qquad e_0=\abs{y_0-1}=1,\quad \eta=\frac1{100}.
\]
\begin{lstlisting}
from fractions import Fraction

q = Fraction(3, 4)        # contraction factor of T(x) = (3x + 1)/4
eta = Fraction(1, 100)    # the bias added at every update
y = Fraction(0)           # the starting point y_0
e0 = abs(y - 1)           # initial error; the fixed point of T is 1
for n in range(21):
    closed = Fraction(26, 25) * (1 - q**n)
    assert y == closed                  # closed form, checked exactly
    error = abs(y - 1)
    bound = q**n * e0 + eta * (1 - q**n) / (1 - q)
    assert error <= bound               # bound, checked exactly
    print(n, round(float(error), 4), round(float(bound), 4))
    y = (3 * y + 1) / 4 + eta           # the biased update gives y_{n+1}
\end{lstlisting}
A line \code{assert} \emph{condition} does nothing when the condition is true, and stops the program with an error message when it is false. So if the program displays all $21$ lines, the closed form and the bound have both been confirmed, in exact arithmetic, for $n=0,1,\ldots,20$. At $n=0$ the inequality reads $1\le1$, which needs no theorem; part~2 of Theorem~\ref{thm:inexact} covers every $n\ge1$.

The function \code{float} converts a fraction to the nearest floating-point number, and \code{round(}\dots\code{, 4)} rounds it to four decimal places. Both happen only in the \code{print} line, after the exact checks, so rounding affects only what is displayed. The displayed bounds are approximations, some of them rounded down; a bound used as a guarantee must be rounded up instead, as Section~\ref{ch12:sec:example} explained. Here is part of the output:
\par\noindent\begin{minipage}{\linewidth}
\begin{lstlisting}
0 1.0 1.0
1 0.74 0.76
2 0.545 0.58
...
10 0.0186 0.0941
11 0.0039 0.0805
12 0.0071 0.0704
...
20 0.0367 0.043
\end{lstlisting}
\end{minipage}\par
Python leaves out trailing zeros, so \code{0.58} stands for $0.5800$ and \code{1.0} for $1.0000$. The row for $n=1$ shows the decimal forms of the exact values $37/50$ and $19/25$ found by hand in the solution to Exercise~\ref{ex:18.6}. As Section~\ref{ch18:sec:experiment} describes, the true error in the middle column falls until $n=11$ and then rises again toward $1/25=0.04$. The bound in the last column, by contrast, decreases at every step toward the same limit.

\section{Checking the Lean examples}\label{appD:sec:lean}
Install Lean, with the editor VS Code and its Lean extension, by following the official instructions~\cite{leaninstall} (Section~\ref{ch15:sec:checker}). The examples are written for Lean version~4.19.0. Lean changes from one version to the next, so if an example that you copied exactly is rejected, check the version first: \code{lean -{}-version} in a terminal shows it. The installation includes a small program called \emph{elan}, which manages the versions of Lean and lets each folder choose its own. To use version~4.19.0, choose a folder for your Lean work and create in it a plain text file named \code{lean-toolchain} whose only line is
\par\noindent\begin{minipage}{\linewidth}
\begin{lstlisting}
leanprover/lean4:v4.19.0
\end{lstlisting}
\end{minipage}\par
From then on, whenever Lean is started in that folder, elan uses version~4.19.0, downloading it first if necessary.

Then create, in the same folder, a file named \code{proofs.lean} with the shape of Section~\ref{ch15:sec:checker}:
\par\noindent\begin{minipage}{\linewidth}
\begin{lstlisting}
import Std

namespace MathCourse

-- the definitions and theorems go here

end MathCourse
\end{lstlisting}
\end{minipage}\par
Everything from \code{-{}-} to the end of a line is a comment, which Lean ignores. Between the two lines that contain \code{MathCourse}, copy the Lean \emph{declarations}, the displayed blocks of code that introduce a definition with \code{def} or a theorem with \code{theorem}, from Chapter~\ref{ch:lean1}, then Chapter~\ref{ch:lean2}, then the solutions for Weeks~15 and~16 in Appendix~\ref{app:solc}, each in the order of the book. The order matters because Lean reads the file from top to bottom and accepts a name only after it has been defined; for example, the definition of \code{Inj} in Section~\ref{ch15:sec:injcomp} must come before \code{inj\_comp} and every later theorem that uses \code{Inj}.

Leave out the other displays: some show what Lean prints, such as a proof state or an error message, and a few repeat lines of a longer proof. The \code{\#print axioms} commands of Section~\ref{ch16:sec:axioms} are discussed below, and the one-line \code{theorem bad} of that section does not belong in the file, since it uses \code{sorry} on purpose.

\subsection*{Running the check}
Open the file in VS Code. Lean checks it while you type, and when the cursor is inside a proof, the infoview beside the file shows the proof state there. To check the whole file at once from a terminal, move to its folder, as in Section~\ref{appD:sec:python}, and type
\par\noindent\begin{minipage}{\linewidth}
\begin{lstlisting}
lean proofs.lean
\end{lstlisting}
\end{minipage}\par
If the command prints nothing, every declaration was accepted, with no error and no warning. An \emph{error} means that the declaration containing it was rejected; a \emph{warning}, that it was accepted but something in it deserves a second look. With all the declarations in the file, every one is accepted, but Lean prints two warnings, each followed by a note on how to switch such warnings off. The first looks like this:
\par\noindent\begin{minipage}{\linewidth}
\begin{lstlisting}
proofs.lean:98:5: warning: unused variable `hf`
note: this linter can be disabled with `set_option linter.unusedVariables false`
\end{lstlisting}
\end{minipage}\par
Every message begins with the file name, the line and the column of the problem, and then says whether it is an error or a warning. Columns are counted from $0$, so column~$5$ is the sixth character of the line, here the start of the name \code{hf}. Your line number may differ from $98$, depending on the blank lines you leave between declarations. The second warning says the same about \code{hg}. Both concern the theorem \code{surj\_comp\_assumed}, and Section~\ref{ch16:sec:irrelevant} explains why it never needs these two hypotheses.

\subsection*{Asking which axioms a proof uses}
To list the axioms on which a theorem depends (Section~\ref{ch16:sec:axioms}), add a line such as
\par\noindent\begin{minipage}{\linewidth}
\begin{lstlisting}
#print axioms MathCourse.surj_comp
\end{lstlisting}
\end{minipage}\par
after the line \code{end MathCourse}; since it stands outside the namespace, it uses the full name, with the prefix \code{MathCourse}. Lean answers
\par\noindent\begin{minipage}{\linewidth}
\begin{lstlisting}
'MathCourse.surj_comp' does not depend on any axioms
\end{lstlisting}
\end{minipage}\par
and from then on \code{lean proofs.lean} prints this line as well. Some theorems, such as \code{not\_surj\_succ} and \code{succ\_surj\_int} of Section~\ref{ch16:sec:debugging}, list the standard axiom \code{propext} instead, through the library lemmas they use; as that section explains, this is no cause for concern.

The command \code{\#check}, which Section~\ref{ch16:sec:debugging} uses to display the exact statement of a library lemma, goes in the file the same way, on a line of its own outside every proof: the line \code{\#check Nat.succ\_ne\_zero}, for example, makes Lean print the statement of that lemma.

\subsection*{What an unfinished proof looks like}
To see for yourself what Section~\ref{ch16:sec:axioms} describes, temporarily replace the last line of the proof of \code{surj\_comp}, \code{rw [ha, hb]}, by the single word \code{sorry}, keeping the same indentation. Lean still accepts the file, but it prints the warning \code{declaration uses 'sorry'} at the name of the theorem, and the \code{\#print axioms} line above now reports
\par\noindent\begin{minipage}{\linewidth}
\begin{lstlisting}
'MathCourse.surj_comp' depends on axioms: [sorryAx]
\end{lstlisting}
\end{minipage}\par
Put the original line back before you rely on the file again: a file that contains \code{sorry} has not proved what it states.

\subsection*{Where to go next}
Once the statements, goals and proof states of Chapters~\ref{ch:lean1} and~\ref{ch:lean2} feel familiar, a natural next step is \emph{Mathematics in Lean}, an interactive textbook with many exercises~\cite{mil}. It is built on Mathlib, the large mathematical library that this book's examples do without, and its own instructions explain how to set it up.

